\section{הרצאה 13: משוואת פואסון בעיגול ומשוואות אינטגרליות}
\date{20.01.2026}

\begin{mainbox}{מידע כללי}
\textbf{תאריך:} 20/01/2026\\
\textbf{נושאים:} השלמת פתרון משוואת פואסון (שיטת הפריסה), מבוא למשוואות אינטגרליות (משוואת \LR{Fredholm}), גרעינים פריקים, גרעיני הפרש וטרנספורם פורייה.
\end{mainbox}

\subsection*{שאלות מנחות}
\begin{enumerate}[label=\textbf{\arabic*.}, rightmargin=2em, leftmargin=1em]
\item כיצד פותרים את משוואת פואסון $\nabla^2 u = f$ כאשר לא ניתן לנחש פתרון פרטי בקלות?
\item מהי משוואת \LR{Fredholm} מסוג שני ומה ההבדל בינה לבין מערכת משוואות לינארית?
\item מהו גרעין פריק (\LR{Separable Kernel}) וכיצד הוא מאפשר פתרון אנליטי מדויק?
\item מהם הערכים העצמיים של משוואה אינטגרלית ומתי לא קיים פתרון יחיד?
\item כיצד פותרים משוואה אינטגרלית בתחום אינסופי עם גרעין תלוי הפרש ($K(x-y)$)?
\end{enumerate}

\subsection{משוואת פואסון בעיגול: השיטה הכללית}

בשיעור הקודם ראינו כיצד לפתור את משוואת פואסון $\nabla^2 u = f(r,\theta)$ על ידי ניחוש פתרון פרטי. כעת נציג את השיטה השיטתית \LR{(Brute Force)} המתאימה לכל פונקציית מקור $f$.

\begin{formulabox}{הבעיה}
נרצה לפתור בעיגול $0 \le r < R$:
\[
\nabla^2 u(r, \theta) = f(r, \theta)
\]
עם תנאי שפה $u(R, \theta) = g(\theta)$.
\end{formulabox}

\subsubsection*{אלגוריתם הפתרון (שיטת הפריסה)}

\begin{steps}
\item \textbf{פריסת המקור \LR{(Source Expansion)}:}
נפרוס את הפונקציה הידועה $f(r, \theta)$ לטור פונקציות עצמיות (במקרה זה, טור פורייה זוויתי):
\[
f(r, \theta) = \alpha_0(r) + \sum_{n=1}^{\infty} \left[ \alpha_n(r) \cos(n\theta) + \beta_n(r) \sin(n\theta) \right]
\]
המקדמים $\alpha_n(r), \beta_n(r)$ מחושבים על ידי אינטגרלים (הטלות):
\[
\alpha_n(r) = \frac{1}{\pi} \int_{0}^{2\pi} f(r, \theta) \cos(n\theta) \, d\theta \quad (\text{עבור } n>0)
\]

\item \textbf{פריסת הפתרון:}
נניח פתרון מהצורה:
\[
u(r, \theta) = A_0(r) + \sum_{n=1}^{\infty} \left[ A_n(r) \cos(n\theta) + B_n(r) \sin(n\theta) \right]
\]
כאשר $A_n(r), B_n(r)$ הן פונקציות לא ידועות.

\item \textbf{הצבה במשוואה והשוואת מקדמים:}
הלפלסיאן בקואורדינטות גליליות הוא $\nabla^2 = \partial_{rr} + \frac{1}{r}\partial_r + \frac{1}{r^2}\partial_{\theta\theta}$. נציב ונקבל משוואות דיפרנציאליות רגילות (ODEs) עבור המקדמים הרדיאליים:
\[
A_n''(r) + \frac{1}{r}A_n'(r) - \frac{n^2}{r^2}A_n(r) = \alpha_n(r)
\]
זוהי משוואת אוילר לא-הומוגנית.

\item \textbf{פתרון ה-ODEs:}
\begin{itemize}
    \item \textbf{עבור $n > 0$:} הפתרון ההומוגני הוא $c_1 r^n + c_2 r^{-n}$. הפתרון הפרטי תלוי ב-$\alpha_n(r)$.
    \item \textbf{עבור $n = 0$:} המשוואה היא $\frac{1}{r}(r A_0')' = \alpha_0(r)$.
\end{itemize}
יש לדרוש שהפתרון לא יתבדר בראשית ($r=0$), ולכן זורקים איברים כמו $r^{-n}$ או $\ln r$ (אלא אם יש מקור נקודתי בראשית).
\end{steps}

\begin{examplebox}{דוגמה: מקור קבוע $f(r, \theta) = 1$}
במקרה זה, כל המקדמים $\alpha_n, \beta_n$ מתאפסים למעט $\alpha_0(r) = 1$.
המשוואה עבור $n=0$:
\[
\frac{1}{r}\frac{d}{dr}\left(r \frac{dA_0}{dr}\right) = 1 \implies r \frac{dA_0}{dr} = \frac{r^2}{2} + C_1 \implies A_0(r) = \frac{r^2}{4} + C_1 \ln r + C_2
\]
לפתרון רגולרי נבחר $C_1 = 0$. את $C_2$ ניתן לבחור לנוחות (למשל כך שיתאפס על השפה) או להשאיר ולהתאים בסוף לתנאי השפה הכוללים.
\end{examplebox}

\section{מבוא למשוואות אינטגרליות}

\subsection{מוטיבציה והגדרות}
עד כה עסקנו בבעיות דיפרנציאליות ("בעיות ישירות"). משוואות אינטגרליות מופיעות לעיתים קרובות בבעיות הפוכות (Inverse Problems), למשל: נתונה צורת מיתר, מצא את הכוח שפעל עליו. כמו כן, ניתן להמיר בעיות שפה במד"ר למשוואות אינטגרליות באמצעות פונקציות גרין.

\begin{definitionbox}{משוואת \LR{Fredholm} מסוג שני}
הצורה הכללית של המשוואה (בחד-ממד):
\[
f(x) = g(x) + \lambda \int_{a}^{b} K(x, y) f(y) \, dy
\]
\begin{itemize}
    \item $f(x)$: הפונקציה הנעלמת.
    \item $g(x)$: האיבר האי-הומוגני (המקור).
    \item $K(x, y)$: ה\textbf{גרעין} (\LR{Kernel}) של המשוואה.
    \item $\lambda$: פרמטר (לעיתים משחק תפקיד של ערך עצמי).
\end{itemize}
אם $g(x) = 0$, המשוואה נקראת \textbf{הומוגנית}.
\end{definitionbox}

\begin{notebox}{אינטואיציה: דיסקרטיזציה}
ניתן לחשוב על האינטגרל כגבול של סכום. אם נחלק את הקטע $[a,b]$ ל-$N$ נקודות, המשוואה הופכת למערכת משוואות ליניאריות:
\[
f(x_i) \approx g(x_i) + \lambda \sum_{j=1}^{N} K(x_i, x_j) f(x_j) \Delta x
\]
זהו אנלוג למשוואה המטריציונית $\vec{f} = \vec{g} + \lambda \mathbf{K} \vec{f}$.
\end{notebox}

\subsection{פתרון עבור גרעין פריק (\LR{Separable Kernel})}

שיטת פתרון אנליטית קיימת כאשר הגרעין ניתן להפרדה למכפלת פונקציות של $x$ ופונקציות של $y$.

\begin{definitionbox}{גרעין פריק}
גרעין נקרא פריק אם ניתן לכתוב אותו כסכום סופי:
\[
K(x, y) = \sum_{i=1}^{N} P_i(x) Q_i(y)
\]
\end{definitionbox}

\subsubsection*{אלגוריתם הפתרון}
נציג את השיטה דרך דוגמה מההרצאה.
\begin{examplebox}{דוגמה: משוואה עם גרעין פריק}
נפתור את המשוואה:
\[
f(x) = \underbrace{1 - 6x^2}_{g(x)} + \lambda \int_{-1}^{1} \underbrace{(x + x^2 y^3)}_{K(x,y)} f(y) \, dy
\]
\textbf{שלב 1: הפרדת הגרעין והוצאת גורמים התלויים ב-$x$}
\[
f(x) = 1 - 6x^2 + \lambda x \int_{-1}^{1} f(y) \, dy + \lambda x^2 \int_{-1}^{1} y^3 f(y) \, dy
\]
האינטגרלים הם \textbf{מספרים קבועים} (לא תלויים ב-$x$). נסמנם ב-$A$ ו-$B$:
\[
A \equiv \int_{-1}^{1} f(y) \, dy, \quad B \equiv \int_{-1}^{1} y^3 f(y) \, dy
\]

\textbf{שלב 2: זיהוי מבנה הפתרון}
מתוך המשוואה, אנו יודעים כעת ש-$f(x)$ חייבת להיות מהצורה:
\[
\boxed{f(x) = 1 - 6x^2 + \lambda A x + \lambda B x^2}
\]

\textbf{שלב 3: בניית מערכת משוואות עבור הקבועים}
נציב את הביטוי ל-$f(y)$ בתוך הגדרות האינטגרלים $A$ ו-$B$:
\begin{align*}
A &= \int_{-1}^{1} \left( 1 - 6y^2 + \lambda A y + \lambda B y^2 \right) dy \\
B &= \int_{-1}^{1} y^3 \left( 1 - 6y^2 + \lambda A y + \lambda B y^2 \right) dy
\end{align*}
נחשב את האינטגרלים (תוך שימוש בסימטריה: אינטגרל על פונקציה אי-זוגית בתחום סימטרי מתאפס):
\begin{align*}
A &= [y - 2y^3]_{-1}^{1} + 0 + \lambda B [\frac{y^3}{3}]_{-1}^{1} = -2 + \frac{2}{3}\lambda B \\
B &= 0 + 0 + \lambda A [\frac{y^5}{5}]_{-1}^{1} + 0 = \frac{2}{5}\lambda A
\end{align*}

\textbf{שלב 4: פתרון המערכת האלגברית}
קיבלנו מערכת ליניארית:
\[
\begin{cases}
A - \frac{2}{3}\lambda B = -2 \\
-\frac{2}{5}\lambda A + B = 0
\end{cases}
\]
נפתור (למשל ע"י הצבה $B = 0.4\lambda A$ במשוואה הראשונה):
\[
A\left(1 - \frac{4}{15}\lambda^2\right) = -2
\]
\end{examplebox}

\begin{theorembox}{קיום ויחידות פתרון}
בדוגמה הנ"ל, הפתרון קיים ויחיד כל עוד הדטרמיננטה אינה מתאפסת:
\[
1 - \frac{4}{15}\lambda^2 \neq 0 \implies \lambda \neq \pm \frac{\sqrt{15}}{2}
\]
עבור ערכים קריטיים אלו של $\lambda$ (הערכים העצמיים), אם המשוואה אי-הומוגנית, בדרך כלל לא יהיה פתרון. אם המשוואה הומוגנית ($g=0$), אלו הערכים היחידים עבורם קיים פתרון לא טריוויאלי.
\end{theorembox}

\section{משוואות עם גרעין הפרש ותחום אינסופי}

סוג נוסף של משוואות הוא כאשר הגרעין תלוי רק בהפרש הקואורדינטות, $K(x,y) = K(x-y)$, והתחום הוא $(-\infty, \infty)$.
המשוואה היא:
\[
f(x) = g(x) + \lambda \int_{-\infty}^{\infty} K(x-y) f(y) \, dy
\]
האינטגרל הוא למעשה \textbf{קונבולוציה} (\LR{Convolution}).

\subsection{שיטת הפתרון: טרנספורם פורייה}

נשתמש בהגדרת הטרנספורם הבאה (לפי הקונבנציה בכיתה):
\[
\tilde{f}(k) = \int_{-\infty}^{\infty} f(x) e^{-ikx} \, dx, \quad f(x) = \frac{1}{2\pi} \int_{-\infty}^{\infty} \tilde{f}(k) e^{ikx} \, dk
\]

\begin{steps}
\item \textbf{הפעלת הטרנספורם:} נבצע טרנספורם פורייה על שני אגפי המשוואה.
\item \textbf{משפט הקונבולוציה:} הטרנספורם של הקונבולוציה הוא מכפלת הטרנספורמים:
\[
\mathcal{F}\left\{ \int K(x-y)f(y)dy \right\} = \tilde{K}(k) \cdot \tilde{f}(k)
\]
\item \textbf{פתרון אלגברי במישור $k$:}
\[
\tilde{f}(k) = \tilde{g}(k) + \lambda \tilde{K}(k) \tilde{f}(k)
\]
נבודד את $\tilde{f}(k)$:
\[
\boxed{\tilde{f}(k) = \frac{\tilde{g}(k)}{1 - \lambda \tilde{K}(k)}}
\]
\item \textbf{טרנספורם הפוך:} נבצע טרנספורם הפוך לקבלת $f(x)$.
\end{steps}

\begin{examplebox}{דוגמה: גרעין גאוסי}
נתונה המשוואה:
\[
f(x) = x e^{-x^2} + \lambda \int_{-\infty}^{\infty} e^{-a(x-y)^2} f(y) \, dy
\]
\textbf{פתרון:}
\begin{enumerate}
    \item נחשב את הטרנספורם של $g(x) = x e^{-x^2}$. נזכור כי $x \to i \partial_k$ במישור פורייה, וכי טרנספורם של גאוסיאן הוא גאוסיאן.
    \item נחשב את הטרנספורם של הגרעין $K(z) = e^{-az^2}$. זהו אינטגרל גאוסי סטנדרטי: $\tilde{K}(k) = \sqrt{\frac{\pi}{a}} e^{-k^2/4a}$.
    \item נציב בנוסחה ונבצע טרנספורם הפוך.
\end{enumerate}
\end{examplebox}

\begin{summarybox}{סיכום שיטות}
\begin{itemize}
    \item \textbf{משוואת פואסון:} אם אין ניחוש קל, פורסים את המקור ואת הפתרון לטורי פונקציות עצמיות ופותרים מערכת אינסופית של מד"רים מופרדים.
    \item \textbf{משוואת אינטגרלית (גרעין פריק):} הופכת למערכת משוואות אלגברית ליניארית סופית.
    \item \textbf{משוואה אינטגרלית (גרעין הפרש):} נפתרת אלגברית במישור התדר באמצעות טרנספורם פורייה.
\end{itemize}
\end{summarybox}









\maketitle

\begin{summarybox}{מנהלות למבחן וסוף סמסטר}
\begin{itemize}
    \item \textbf{מבנה המבחן:} ללא הפתעות, דומה לשנים קודמות.
    \item \textbf{חומר עזר:} חומר פתוח. המטרה אינה שינון נוסחאות אלא הבנה. מומלץ להכין דפי סיכום מסודרים. אין הגבלה על כמות הדפים.
    \item \textbf{לו"ז:} נותרו שני שיעורים. נושא אחרון: משוואות אינטגרליות (חומר הנחשב קל יחסית).
    \item \textbf{תרגול:} יפורסם תרגיל רשות בנושא משוואות אינטגרליות עם פתרון לתרגול עצמי.
\end{itemize}
\end{summarybox}

\section{משוואת פואסון בעיגול - שיטה כללית}

בשיעור הקודם עסקנו בפתרון משוואת פואסון על ידי ניחוש פתרון פרטי. כעת נציג את הדרך הכללית (שיטת הפריסה) למקרים בהם לא ניתן לנחש בקלות.

\begin{mainbox}{הבעיה המתמטית}
אנו רוצים לפתור את משוואת פואסון בקואורדינטות קוטביות בתוך עיגול ($0 \le r \le R$):
$$ \nabla^2 u(r, \theta) = f(r, \theta) $$
במקרה הספציפי הנדון:
\begin{itemize}
    \item צפיפות המטען קבועה: $f(r, \theta) = 1$.
    \item רדיוס העיגול: $R=1$.
    \item תנאי שפה: $u(1, \theta) = \cos(3\theta)$.
\end{itemize}
\end{mainbox}

\begin{examplebox}{פתרון בשיטת הפריסה (טורים)}
הרעיון הוא לפרוס את הפונקציה הידועה $f(r, \theta)$ ואת הפתרון המבוקש $u(r, \theta)$ לטור פונקציות עצמיות (במקרה זה, טור פורייה בזווית $\theta$).

\textbf{שלב א': פריסת $f(r, \theta)$}
$$ f(r, \theta) = \alpha_0(r) + \sum_{n=1}^{\infty} \left( \alpha_n(r) \cos(n\theta) + \beta_n(r) \sin(n\theta) \right) $$
במקרה שלנו $f=1$, ולכן כל המקדמים מתאפסים פרט ל-$\alpha_0$:
$$ \alpha_0(r) = 1, \quad \alpha_n(r) = 0 \quad (n>0) $$

\textbf{שלב ב': פריסת הפתרון $u(r, \theta)$}
$$ u(r, \theta) = A_0(r) + \sum_{n=1}^{\infty} \left( A_n(r) \cos(n\theta) + B_n(r) \sin(n\theta) \right) $$

\textbf{שלב ג': הצבה במשוואה והשוואת מקדמים}
משוואת פואסון בקואורדינטות פולריות:
$$ u_{rr} + \frac{1}{r}u_r + \frac{1}{r^2}u_{\theta\theta} = f $$
לאחר הצבת הטורים וגזירה איבר-איבר, נקבל משוואות דיפרנציאליות רגילות עבור המקדמים $A_n(r)$:
$$ A_n''(r) + \frac{1}{r}A_n'(r) - \frac{n^2}{r^2}A_n(r) = \alpha_n(r) $$
(וכנ"ל עבור $B_n$).

\textbf{פתרון עבור $n > 0$:}
אגף ימין ($\alpha_n$) הוא 0. זו משוואת אוילר הומוגנית.
הפתרון הכללי: $A_n(r) = c_1 r^n + c_2 r^{-n}$.
משיקולי אי-התבדרות בראשית ($r=0$), נזרוק את $r^{-n}$.
$$ A_n(r) = c_n r^n $$

\textbf{פתרון עבור $n = 0$ (האיבר האי-הומוגני):}
המשוואה היא:
$$ A_0''(r) + \frac{1}{r}A_0'(r) = 1 $$
נכפיל ב-$r$ ונקבל: $(r A_0')' = r$. נבצע אינטגרציה כפולה:
1. $r A_0' = \frac{r^2}{2} + C_1$
2. $A_0' = \frac{r}{2} + \frac{C_1}{r}$
3. $A_0(r) = \frac{r^2}{4} + C_1 \ln r + C_2$

משיקולי אי-התבדרות, $C_1=0$.
קיבלנו: $A_0(r) = \frac{r^2}{4} + C_2$.
ניתן לבחור את $C_2$ כרצוננו (משפיע על תנאי השפה של החלק ההומוגני). נבחר $C_2 = -1/4$ כדי לאפס את החלק הפרטי על השפה $r=1$.

\textbf{שלב ד': התאמה לתנאי השפה}
הפתרון הכללי שקיבלנו:
$$ u(r, \theta) = \left(\frac{r^2}{4} - \frac{1}{4}\right) + \sum_{n=1}^{\infty} A_n r^n \cos(n\theta) $$
תנאי השפה ב-$r=1$ הוא $u(1, \theta) = \cos(3\theta)$.
נציב $r=1$:
$$ 0 + \sum A_n (1)^n \cos(n\theta) = \cos(3\theta) $$
מכאן נסיק שרק המקדם $A_3$ שונה מאפס ושווה ל-1.
$$ A_3 = 1, \quad A_n = 0 \quad \forall n \neq 3 $$

\textbf{תוצאה סופית:}
$$ u(r, \theta) = \left(\frac{r^2}{4} - \frac{1}{4}\right) + r^3 \cos(3\theta) $$
\end{examplebox}

\begin{physicsnote}{השוואה בין השיטות}
\begin{itemize}
    \item \textbf{ניחוש:} מהיר יותר, דורש אינטואיציה.
    \item \textbf{פריסה (\LR{Brute Force}):} תמיד עובדת, שיטתית, אך ארוכה יותר ודורשת פתרון משוואות דיפרנציאליות וחישוב אינטגרלים למקדמים.
\end{itemize}
\end{physicsnote}

\section{משוואות אינטגרליות}

זהו הנושא האחרון בקורס. המוטיבציה מגיעה מ"בעיות הפוכות": אם במשוואה דיפרנציאלית אנו מחפשים את הצורה $y(x)$ בהינתן הכוח $f(x)$, במשוואה אינטגרלית לעיתים נרצה למצוא את הכוח $f(x)$ שגרם לצורה נתונה.

\begin{definitionbox}{משוואת פרדהולם (\LR{Fredholm}) לינארית}
הצורה הכללית של המשוואה:
$$ h(x)f(x) = g(x) + \lambda \int_{a}^{b} K(x,y)f(y) dy $$
\begin{itemize}
    \item $f(x)$: הפונקציה הנעלמת.
    \item $K(x,y)$: הגרעין (\LR{Kernel}) של המשוואה.
    \item $\lambda$: פרמטר (לעיתים משחק תפקיד של ערך עצמי).
    \item אם $h(x)=0$: משוואה מסוג ראשון.
    \item אם $h(x)=1$: משוואה מסוג שני (הנפוצה יותר).
    \item אם $g(x)=0$: המשוואה הומוגנית.
\end{itemize}
\end{definitionbox}

\subsection{גרעין פריק (\LR{Separable Kernel})}
אם הגרעין ניתן לכתיבה כסכום סופי של מכפלות פונקציות של $x$ ו-$y$:
$$ K(x,y) = \sum_{i=1}^{N} M_i(x) N_i(y) $$
אזי המשוואה פתירה אנליטית על ידי רדוקציה למערכת משוואות אלגבריות.

\begin{examplebox}{דוגמה: משוואה אי-הומוגנית עם גרעין פריק}
$$ f(x) = \lambda \int_{-1}^{1} (xy + x^2 y^3)f(y) dy + (1 - 6x^2) $$
זיהוי הגרעין: $K(x,y) = x \cdot y + x^2 \cdot y^3$. זהו גרעין פריק עם $N=2$.

\textbf{שלב 1: הוצאת תלות ב-$x$ מחוץ לאינטגרל}
$$ f(x) = \lambda x \underbrace{\int_{-1}^{1} y f(y) dy}_{A} + \lambda x^2 \underbrace{\int_{-1}^{1} y^3 f(y) dy}_{B} + 1 - 6x^2 $$
האינטגרלים המסומנים הם \textbf{קבועים} (מספרים), נסמנם ב-$A$ ו-$B$.
קיבלנו את התבנית של הפתרון:
$$ f(x) = \lambda A x + (\lambda B - 6)x^2 + 1 $$

\textbf{שלב 2: בניית מערכת משוואות עבור הקבועים $A, B$}
נציב את הביטוי של $f(x)$ בחזרה להגדרות של $A$ ו-$B$:
1. $A = \int_{-1}^{1} y f(y) dy = \int_{-1}^{1} y [\lambda A y + (\lambda B - 6)y^2 + 1] dy$
2. $B = \int_{-1}^{1} y^3 f(y) dy = \int_{-1}^{1} y^3 [\lambda A y + (\lambda B - 6)y^2 + 1] dy$

לאחר חישוב האינטגרלים (איברים אי-זוגיים מתאפסים בתחום סימטרי):
$$ A = \int_{-1}^{1} \lambda A y^2 dy = \lambda A \left[\frac{y^3}{3}\right]_{-1}^{1} = \frac{2}{3}\lambda A \quad \text{(תיקון מהתמליל: החישוב בתמליל הניב תוצאות שונות במעט, נצמד לעקרון)} $$
\textit{הערה: בהרצאה בוצע חישוב ספציפי עבור הדוגמה הנתונה שנתן:}
משוואה 1: $A - 2\lambda B = 0$
משוואה 2: $-\frac{2}{7}\lambda A + B = 0$

\textbf{שלב 3: פתרון המערכת}
זו מערכת לינארית עבור $A, B$. הדטרמיננטה קובעת אם קיים פתרון יחיד.
$$ \det \begin{pmatrix} 1 & -2\lambda \\ -\frac{2}{7}\lambda & 1 \end{pmatrix} = 1 - \frac{4\lambda^2}{7} $$
\begin{itemize}
    \item אם הדטרמיננטה $\neq 0$: קיים פתרון יחיד עבור $A, B$ (במקרה הומוגני הפתרון הוא הטריוויאלי 0, במקרה אי-הומוגני מוצאים את הערכים).
    \item אם הדטרמיננטה $= 0$: אין פתרון או אינסוף פתרונות. הערכים המאפסים הם $\lambda = \pm \sqrt{7}/2$.
\end{itemize}
\end{examplebox}

\subsection{משוואה הומוגנית וערכים עצמיים}
אם $g(x)=0$, המשוואה היא:
$$ f(x) = \lambda \int K(x,y)f(y) dy $$
במקרה זה, למערכת האלגברית שתתקבל יהיה פתרון לא טריוויאלי ($f(x) \neq 0$) \textbf{רק} עבור ערכי $\lambda$ המאפסים את הדטרמיננטה. אלו הם הערכים העצמיים של המשוואה האינטגרלית. הפתרון נקבע עד כדי קבוע כפלי.

\section{משוואות עם גרעין הפרש (\LR{Difference Kernel}) בתחום אינסופי}

סוג נוסף של משוואות הוא כאשר הגרעין תלוי רק בהפרש הקואורדינטות, $K(x,y) = K(x-y)$, והתחום הוא $(-\infty, \infty)$. משוואות אלו נפוצות בפיזיקה (למשל, כוחות התלויים במרחק).

\begin{mainbox}{השיטה: התמרת פורייה}
המשוואה:
$$ f(x) = g(x) + \lambda \int_{-\infty}^{\infty} K(x-y)f(y) dy $$
האינטגרל הוא למעשה \textbf{קונבולוציה} (\LR{Convolution}) בין $K$ ל-$f$.

\textbf{שלבי הפתרון:}
1. נבצע התמרת פורייה על שני אגפי המשוואה. נגדיר: $\tilde{f}(k) = \int_{-\infty}^{\infty} f(x) e^{-ikx} dx$.
2. לפי משפט הקונבולוציה, התמרת הקונבולוציה היא מכפלת ההתמרות:
$$ \tilde{f}(k) = \tilde{g}(k) + \lambda \tilde{K}(k) \cdot \tilde{f}(k) $$
3. נבודד את $\tilde{f}(k)$ בצורה אלגברית:
$$ \tilde{f}(k) = \frac{\tilde{g}(k)}{1 - \lambda \tilde{K}(k)} $$
4. נבצע התמרת פורייה הפוכה כדי למצוא את $f(x)$.
\end{mainbox}

\begin{examplebox}{דוגמה: גרעין גאוסי}
$$ f(x) = x e^{-x^2} + \lambda \int_{-\infty}^{\infty} e^{-a(x-y)^2} f(y) dy $$
\begin{enumerate}
    \item נחשב את $\tilde{g}(k)$ עבור $g(x) = x e^{-x^2}$. (דורש אינטגרציה עם השלמה לריבוע). התוצאה פרופורציונלית ל-$k e^{-k^2/4}$.
    \item נחשב את $\tilde{K}(k)$ עבור $K(z) = e^{-az^2}$. התוצאה היא גאוסיאן במרחב התדר: $\sqrt{\frac{\pi}{a}} e^{-k^2/(4a)}$.
    \item נציב בנוסחה הסופית עבור $\tilde{f}(k)$.
    \item (בשלב זה בדרך כלל משאירים את התשובה כביטוי במרחב התדר או מבצעים אינטגרציה נומרית/אנליטית אם זה פשוט).
\end{enumerate}
\end{examplebox}

\begin{summarybox}{סיכום השיעור הבא}
בשיעור הבא נלמד על:
\begin{enumerate}
    \item הכללה לגבולות אינטגרציה התלויים ב-$x$.
    \item תורת ההפרעות עבור $\lambda$ קטן (כאשר הגרעין אינו פריק).
\end{enumerate}
\end{summarybox}